% from ega0-20-fr.tex
\oldpage[0\textsubscript{IV-1}]{125}
par celle de $B\otimes_A B$ est moins fine que la topologie sur $\mathfrak{J}$
déduite de celle de $B$. Pour démontrer la seconde assertion, écrivons
$M\otimes N$, par abus de notation, le sous-module
$\operatorname{Im}(M\otimes_A N)$ pour deux sous-$A$-modules $M$, $N$ de $B$.
Utilisant la relation
\[
\begin{aligned}
  (xy)\otimes z-x\otimes(yz)
  &=(x\otimes1)(1\otimes z)(y\otimes1-1\otimes y)\\
  &=xz.(y\otimes1-1\otimes y)
    -x.(z\otimes1-1\otimes z)(y\otimes1-1\otimes y)
\end{aligned}
\]
dans le $B$-module $B\otimes_A B$ (défini par $j_1$), on voit aussitôt,
compte tenu de \sref{0.20.4.4-fr}, que, si $\mathfrak{R}$ est un idéal de
$B$, on a
\[
  ((\mathfrak{R}^2\otimes B)+(B\otimes\mathfrak{R}^2))\cap\mathfrak{J}
  \subset(\mathfrak{R}\otimes\mathfrak{R})\cap\mathfrak{J}
  +\mathfrak{R}\mathfrak{J}+\mathfrak{J}^2
\]
et d'autre part, si $x_i$, $y_i$ sont des éléments de $\mathfrak{R}$ tels
que $\sum_i(x_i\otimes y_i)\in\mathfrak{J}$, il résulte de
\sref{0.20.4.4.1-fr} que l'on a
$\sum_i(x_i\otimes y_i)\in\mathfrak{R}\mathfrak{J}$, si bien que finalement
\begin{equation}
\tag{20.4.5.1}
\label{0.20.4.5.1-fr}
  (\mathfrak{R}^2\otimes B+B\otimes\mathfrak{R}^2)\cap\mathfrak{J}
  \subset\mathfrak{R}\mathfrak{J}+\mathfrak{J}^2.
\end{equation}
Or, on a un système fondamental de voisinages de $0$ dans $\mathfrak{J}$
(pour la topologie induite par celle de $B\otimes_A B$) en prenant pour
voisinages de $0$ les ensembles
$(\mathfrak{R}\otimes B+B\otimes\mathfrak{R})\cap\mathfrak{J}$, où
$\mathfrak{R}$ parcourt l'ensemble des idéaux ouverts. Comme la topologie de
$\Omega^1_{B/A}$ déduite de celle de $B$ est aussi le quotient par
$\mathfrak{J}^2$ de la topologie de $\mathfrak{J}$ déduite de celle de $B$,
l'hypothèse sur les idéaux ouverts de $B$ et la relation
\sref{0.20.4.5.1-fr} achèvent de prouver la seconde assertion.
\end{proof}

\begin{definition}[20.4.6]
\label{0.20.4.6-fr}
Soient $p_1$ et $p_2$ les applications composées
$B\xrightarrow{j_1}B\otimes_A B\to\mathrm{P}^1_{B/A}$,
$B\xrightarrow{j_2}B\otimes_A B\to\mathrm{P}^1_{B/A}$, qui sont des
$A$-homomorphismes continus tels que
$\varepsilon\circ p_1=\varepsilon\circ p_2=\mathrm{I}_B$. Le
$A$-homomorphisme continu de $A$-modules
\begin{equation}
\tag{20.4.6.1}
\label{0.20.4.6.1-fr}
  d_{B/A}=p_1-p_2:B\to\Omega^1_{B/A}
\end{equation}
est appelé la \emph{différentielle extérieure de $B$ relative à $A$} ; pour
tout $x\in B$, $d_{B/A}(x)$ (aussi noté $d(x)$ ou $dx$) est appelé la
\emph{différentielle de $x$} (relativement à $A$).
\end{definition}

Lorsque $A=\mathbf{Z}$, on écrit $d_B$ au lieu de $d_{B/\mathbf{Z}}$ ; si $B$
est une algèbre sur un corps premier $P$, on a $d_B=d_{B/P}$.

\begin{proposition}[20.4.7]
\label{0.20.4.7-fr}
Le $B$-module $\Omega^1_{B/A}$ est engendré par les éléments $d_{B/A}(x)$,
où $x$ parcourt un système de générateurs de la $A$-algèbre $B$.
\end{proposition}

\begin{proof}
Comme $d_{B/A}(x)$ est l'image canonique de $x\otimes1-1\otimes x$ dans
$\Omega^1_{B/A}$, la proposition est une conséquence immédiate de
\sref{0.20.4.4-fr}.
\end{proof}

\begin{theorem}[20.4.8]
\label{0.20.4.8-fr}
Soient $A$ un anneau topologique, $B$ une $A$-algèbre topologique.
\begin{enumerate}
  \item[\rm(i)] Il existe un isomorphisme unique de $B$-algèbres augmentées
  topologiques
\begin{equation}
\tag{20.4.8.1}
\label{0.20.4.8.1-fr}
  \varphi:\mathrm{P}^1_{B/A}\xrightarrow{\sim}D_B(\Omega^1_{B/A})
\end{equation}
qui se réduit à l'identité dans $\Omega^1_{B/A}$.
  \item[\rm(ii)] L'homomorphisme $d_{B/A}$ est une $A$-dérivation de $B$
  dans $\Omega^1_{B/A}$, ayant la propriété universelle suivante : pour tout
  $B$-module topologique $L$, l'application $u\mapsto u\circ d_{B/A}$ est
  un isomorphisme de $A$-modules
\begin{equation}
\tag{20.4.8.2}
\label{0.20.4.8.2-fr}
  \operatorname{Hom.\ cont}_B(\Omega^1_{B/A},L)
  \xrightarrow{\sim}\operatorname{D\acute{e}r.\ cont}_A(B,L).
\end{equation}
\end{enumerate}
\end{theorem}

\begin{proof}
